\section{Discussion}

\noindent Methylation clocks are promising biomarkers of aging and social stress, and as tools for mechanistic studies of diseases related to the aging process. Despite their widespread use in these applications, methylation clocks have not been comprehensively evaluated in diverse human groups. These groups are underrepresented in genetic and genomic databases and underserved in biomedicine in terms of access to and quality of healthcare. 

In this study, we sought to evaluate the performance of commonly used methylation clocks in genetically admixed individuals from the Americas in the MAGENTA AD study. We found that most clocks did not predict age as accurately in admixed individuals as in White individuals, especially for cohorts with substantial African genetic ancestry. These results replicated in three external, independent cohorts: one of White individuals from Sweden and two of African Americans. We next found that most methylation clocks could not consistently distinguish AD patients and non-demented controls based on age acceleration metrics in non-White cohorts with admixed genetic ancestries.

To evaluate potential genetic factors that could contribute to this decrease in performance, we hypothesized that two types of variants could reduce clock accuracy: 1) variants that disrupt clock CpG sites and prevent methylation and 2) meQTLs that influence clock CpG site methylation. Both scenarios could lead to over or under estimates of age, and therefore spurious associations with age-related disease, if they differ in frequency across genetic ancestries. We discovered that 245 of the 353 CpG sites used by the Horvath clock are disrupted in at least one individual in gnomAD, but these variants are extremely rare and thus unlikely to be a major driver of differences between the cohorts. In contrast, thousands of meQTLs from multiple global populations affected the Horvath clock CpG sites. Many of these meQTL are common, and they are more frequent in individuals with African ancestry in both gnomAD and MAGENTA. We also found that many clock CpGs affected by meQTL also have methylation levels that are significantly associated with clock error in the MAGENTA cohorts. However, we note that the presence of known meQTL was not sufficient to explain all observed clock error at the individual level.
%We also showed that these variants, as well as many that are specific to African ancestry individuals, are present in the MAGENTA study individuals and influence methylation levels in a way that is consistent with their reported betas. Importantly, .

% We also showed that they are  most frequent in African local ancestry blocks in admixed Latino individuals from the Americas.

Our findings demonstrate that methylation clocks---a widespread tool in aging, genomics, and social epidemiology research---perform inconsistently across individuals of different genetic ancestries. These results further underline the need for more genetic diversity in the development and evaluation of genomic and epigenomic tools for precision medicine.

We hope that these results also encourage researchers using these tools to exercise caution when interpreting differences in age acceleration. We have shown that many methylation clocks differ significantly in their accuracy at predicting age between cohorts. Thus, what might appear to be a faster pace of aging, could simply be the result of a difference in genetic ancestry from the training cohort. Broadly applying existing methylation clocks could lead to grave consequences and exacerbate existing disparities in access to quality healthcare, as well as provide spurious conclusions about an individual's health. Due to the increased potential for false positives and false negatives when applying the clocks as predictive biomarkers, individuals at risk might not receive the medical attention they need, and additional stress could unncessarily be placed on individuals in good health. These challenges must be addressed before methylation clocks are adopted as biomarkers for precision medicine. 

The challenges we identify here for methylation clocks mirror the limitations of PRS, wherein phenotype prediction models decrease in accuracy on individuals genetically distant from the training population. While there are substantial biological differences in the processes modeled by methylation clocks and PRS, we are optimistic that recent progress on building PRS that are more portable across cohorts will provide strategies for improving methylation clocks. 

Our results suggest two promising approaches for building more robust clocks. First, we encourage including individuals from multiple genetic ancestries in the training cohorts. The ability of the PhenoAge clock, which included African Americans in the training cohort, to detect significant age acceleration in the African American AD cases suggests this may improve generalizability. Second, given the large number of meQTL in the human genome and their differences in frequency across human populations, training clocks only on CpG sites that do not have known or population-differentiated meQTL may yield more generalizable clocks. Given that the biological signatures driving methylation clock performance appear to be detectable over large fractions of the genome, removing CpGs with strong meQTL may not substantially limit overall performance. Supporting this, a methylation-based age predictor that is minimally influenced by meQTL performed well in African individuals \citep{meeks_common_2025}. Moreover, the strong and relatively consistent performance across cohorts of the DunedinPACE clock, which lacks CpGs with meQTL, is also consistent with this approach. It also suggests that methylation clocks that predict the pace of aging (rather than age itself) may be more robust, but further work is needed to validate this hypothesis.  

We explored two additional methods for improving generalizability across cohorts that did not yield promising results: ensemble methylation clocks and PC-based methylation clocks. First, motivated by the success of ensembling PRSs for increased accuracy and portability \citep{truong_integrative_2024}, we tested whether combining values from multiple methylation clocks could improve risk prediction. For example, combining the predictions of a biomarker-based clock and a first-generation clock could integrate strengths of multiple biomarkers associated with mortality (e.g., PhenoAge) in addition to the chronological aging process and its biological underpinnings (e.g. Horvath clock). However, this approach did not yield improved predictions, though other ensembling approaches could produce different results. We also tested the PC versions of multiple clocks, including the Horvath clock, owing to their reduction in technical noise and variability \citep{higgins-chen_computational_2022}. However, these clock did not increase accuracy relative to their non-PC versions in the admixed MAGENTA cohorts, and in some cases they performed worse. %The PC versions of the clocks are also based on the underlying selected CpG sites, which are subject to differences in meQTL across populations. Thus, this result further supports the approach of building generalizable methylation clocks through training on diverse, global populations and focusing on methylation sites without known meQTL.

There are several caveats and limitations to our study that we hope future work will address. First, the impact of environment on methylation clock accuracy is unresolved. Differences in environmental factors for both local and global populations might lead to decreases in methylation clock accuracy, but they are difficult to study with the data available. Given this, we focused on genetic influences on CpG sites that could lead to spurious associations between populations. The genetic factors we investigated are not sufficient alone to explain differences in clock performance across populations, more work is needed to investigate other non-genetic factors that might cause methylation clocks to not generalize across individuals. Many of the CpGs we discovered with methylation levels associated with clock error do not have known meQTL or variants. We also note that the methods for accounting for cell type composition heterogeneity in blood have not been extensively evalauted across individuals from different populations. Second, while we attempted to evaluate a representative set of first-, second-, and third-generation clocks, we were not able to evaluate all methylation clocks. In particular, some with closed source code that were only available as a web server could not be used due to data privacy restrictions for the MAGENTA samples. Another limitation is the use of blood samples to generate methylation data for the study of a neurodegenerative disease focused on the central nervous system. However, we note that all methylation clocks tested in the present study were developed using blood samples, either exclusively (Hannum, PhenoAge, Zheng\_EN, Zheng\_BLUP, and DunedinPACE) or as part of the tissues used in training (Horvath). In addition, these clocks and their association with age-related diseases such as AD have all been validated in multiple tissues, including blood. Multiple studies in the AD literature point to signatures in blood such as gene expression changes, immune cell type composition, and disruption of the blood brain barrier in AD patients relative to non-demented controls, such that signals related to AD pathology can be identified from this peripheral tissue \citep{shigemizu_identification_2022, griswold_immune_2020} and through methylation age acceleration \citep{marioni_dna_2015, raina_cerebral_2017, hodgson_epigenetic_2017}.

Finally, while the MAGENTA study is an excellent resource for exploring methylation clocks and AD in admixed individuals, it is not representative of all genetic ancestries and combinations. Moreover, the sample size of the MAGENTA study is relatively small. Nonetheless, previous studies have found associations between age acceleration and AD in similar sized cohorts. Both PhenoAge \citep{levine_epigenetic_2018} and the Horvath clock \citep{levine_epigenetic_2015} identified age acceleration of less than a year in AD patients relative to non-demented individuals in cohorts of similar sample sizes (e.g., 700 for Levine et al. 2015 and 604 for Levine et al. 2018). We replicated the association with AD in the white MAGENTA cohort, demonstrating the ability to detect methylation effects in MAGENTA. We also performed power calculations based on an effect of the size observed in previous studies (0.5 year acceleration). Stratifying by MAGENTA cohorts, we had 75\% power for the African Americans, 72\% for the Puerto Ricans, 72\% for the Whites, 65\% for the Peruvians, and 47\% for the Cubans. Thus, it would have been very unlikely to observe no significant differences in any of the admixed cohorts if the effect sizes were similar to previous studies. Thus, reduced clock accuracy is likely a contributor to the lack of association. Furthermore, recent complementary findings on the decreased performance of methylation clocks at age prediction and the role of meQTL in African cohorts \citep{meeks_common_2025} further support our conclusions.

%Moreover, the sample sizes vary between the MAGENTA cohorts, with the Cuban and Peruvian cohorts being particularly small relative to the others.

In conclusion, our results show that many existing methylation clocks have inconsistent performance and limited portability across genetically admixed cohorts. We encourage future efforts in the development of methylation clocks and other genomics-based aging biomarkers to be genetics- and ancestry-aware to ensure the accuracy of these tools for all individuals, regardless of their genetic background.

\clearpage